\input{preamble.tex}

% Title info
\title{Transmittance Coefficient Generation Package\\ \normalsize{Quickstart Guide}}
\date{\today}
\docnumber{(unassigned)}
\docseries{CRTM}


%-------------------------------------------------------------------------------
%                            Ze document begins...
%-------------------------------------------------------------------------------
\begin{document}

\maketitle

%\draftwatermark

\tableofcontents

\newpage

% The front matter
%=================
\thispagestyle{empty}
\vspace*{10cm}
\begin{center}
  {\sffamily\Large\bfseries Change History}
  \begin{table}[htp]
    \centering
    \begin{tabular}{l l l}
      \hline
      \sffamily\textbf{Date} & \sffamily\textbf{Author} & \sffamily\textbf{Change}\\
      \hline\hline
      2021-12-14 & P. Stegmann& Initial Draft.\\
      2024-06-07 & B. Johnson & (bjohns@ucar.edu or Benjamin.T.Johnson@noaa.gov) \\
       & & Updated for improved workflow.\\
      \hline
    \end{tabular}
  \end{table}
\end{center}
\clearpage
\pagenumbering{arabic}
\setcounter{page}{1}


% The main matter
%================
%\include{}

\section{Introduction}
This document is a quick start guide for the JCSDA CRTM Transmittance Coefficient Generation package (CGP).
To this end, the document provides a minimal step-by-step guide to create CRTM \verb|SpcCoeff| and \verb|TauCoeff| files for an example MW instrument. 
This is a minimal walkthrough that only demonstrates the mechanics and necessary commands for the end-to-end workflow to obtain instrument coefficients.

\section{CRTM Microwave Coefficient Generation}
The example sensor chosen for this walkthrough is TROPICS Pathfinder, with the CRTM instrument ID \emph{tropics\_sv1\_srf\_v2}.

\subsection{Prerequisites}
The CGP requires a number of depencies that need to be available:
\begin{itemize}
  \item A Fortran compiler, e.g. gfortran or ifort.
  \item netCDF4 and HDF5.
  \item Some sort of BLAS package, such as openblas (brew install openblas, also provided by HPC module systems (e.g., spack-stack, hpc-stack, etc.)
  \item The slurm scheduler for the ODPS regression.
  \item The cmake build system.
  \item git
\end{itemize}

\subsection{Compiling the code}
Before running any of the Fortran applications, the above items need to be compiled and their dependencies enabled.

\subsubsection{Cloning the Coefficient Generation Package}
The CGP can be cloned like this:
\begin{verbatim}
 git clone https://github.com/JCSDA-internal/CRTM_coef
\end{verbatim}

\subsubsection{Creating a new git branch}
Before you create a new instrument coefficient, it is highly encouraged that you create a new branch of the \verb|CRTM_coef| repository for your work, e.g.:
\begin{verbatim}
 git branch feature/tropics
 git checkout feature/tropics
\end{verbatim}

\subsubsection{Configuring the Coefficient Generation Package}
Currently the build system for the CGP is based on \verb|cmake|.  It will attempt to find BLAS and netCDF.

\begin{verbatim}
 # Entering the CGP root directory and starting the \verb|cmake| process.
 cd CRTM_coef/
 mkdir build
 cd build
 cmake ..
\end{verbatim}

You might run into issues where cmake can't find netCDF.
\begin{verbatim}
 #help cmake find netcdf
 # example:
 > locate netcdf.mod
 /opt/netcdf4/4.4.1-intel-16.0.2/include/netcdf.mod
 /opt/netcdf4/4.6.2-intel-16.0.2/include/netcdf.mod
 /opt/netcdf4/4.6.2-intel-17.0.6/include/netcdf.mod
 /opt/netcdf4/4.6.2-intel-18.0.3/include/netcdf.mod
 /opt/netcdf4/4.7.1-intel-19.0.5/include/netcdf.mod
 /opt/netcdf4/4.7.3-intel-18.0.3/include/netcdf.mod
 /opt/netcdf4/4.7.3-intel-19.0.5/include/netcdf.mod
 /opt/netcdf4/4.8.1-intel-2022.1/include/netcdf.mod

 #pick the path to the include directory that matches your compiler 
 export NetCDF_INCLUDE_DIRS=/opt/netcdf4/4.7.1-intel-19.0.5/include
\end{verbatim}
 


Aside for macos/linux box (tested on macbook-intel so far).  This package is really designed for HPC use, because it's fairly intensive.  However, for MW-only development, you can probably get away with doing it on a Mac/Linux desktop / laptop. A lot of the instruction / testing has occured with Intel Fortran Compiler (ifort), but you can give this a shot on your local box
\begin{verbatim}
 brew install gcc
 brew install netcdf
 brew install netcdf-fortran
 brew install openblas
\end{verbatim}
OR install spack-stack:  https://wiki.ucar.edu/display/JEDI/Building+spack-stack+on+MacOS  (this takes awhile, and is unnecessary if you're only doing CRTM stuff).  

Try to run the \verb|cmake ..| step again.

Contact the JCSDA CRTM team via github with any questions (\verb|https://github.com/JCSDA/CRTM_coef|)

\begin{verbatim}
 #still in the \verb|build| directory ...
 make -j12
 # if successful
 make install
\end{verbatim}

Have a look at the \verb|build/bin| directory, you should see the following applications installed:

\subsubsection{Compiling the CGP Fortran Applications}
This is a checklist of applications that are created in the previous step:
\begin{todolist}

\item \verb|Assemble_ODAS|  (occasionally called ``\verb|Cat_ODAS|'' in older documentation)
\item \verb|Assemble_ODPS|
\item \verb|AtmProfile_CreateFile|
\item \verb|Create_SpcCoeff|
\item \verb|GetSenInfo|
\item \verb|MW_TauProfile|
\item \verb|ODAS_NC2BIN|
\item \verb|ODAS_Regress|
\item \verb|ODAS_WLO_Regress|
\item \verb|ODPS_NC2BIN|
\item \verb|ODPS_Regress|
\item \verb|SpcCoeff_NC2BIN|
\item \verb|oSRF_Create_from_ASCII| (occasionally called ``\verb|var.out|'' in older documentation)

\end{todolist}


\subsection{Preparing the SRF}
Before starting the coefficient generation process, the spectral response function (SRF) of the instrument needs to be brought into the correct input format. 
This is the responsibility of the user, as is making sure that the SRF spectral resolution is high enough to resolve all desired absorption features.
Input data for a new CRTM instrument should be stored in the \verb|staging/fix/config/| directory of the CGP.
An example ASCII input file for TROPICS (e.g., \verb|staging/fix/config/TROPICS_Pathfinder/tropics_sv1_srf_v4/SRFDATA/tropics_sv1_srf_v4_ch1.txt| looks like this:)
\begin{verbatim}
89.484  0.0022256
89.522  0.0026217
89.561  0.0033972
89.599  0.004763
89.638  0.0067661
89.676  0.0095845
89.714  0.01315
89.753  0.016683
89.791  0.018966
89.83   0.020105
89.868  0.02048
89.906  0.019944
89.945  0.0187
89.983  0.017721
90.022  0.019286
90.06   0.024162
90.098  0.028599
90.137  0.031213
90.175  0.031505
90.214  0.029475
...
\end{verbatim}
The first column contains the frequency in units of [GHz] and the second column contains the unitless normalized SRF value.
Recommend perusing the \verb|SRFDATA/| directory noted above just to understand how things are organized.   The creation of these files typically come from you, the coefficient developer, based on whatever SRF information you can obtain from the instrument team for a given instrument.  Your best bet for support at this point is to (a) ask the CRTM teams (JCSDA, STAR), and (b) the instrument science team for the given sensor -- they likely have this information.    When you don't know the SRF, you assume a ``boxcar'', which looks like this:
\begin{verbatim}
84.09999999899999 0.0
84.1 1.0
86.9 1.0
86.90000000100001 0.0
\end{verbatim}
Note the sideband handling for this ``85.5'' GHz channel (+/- 1.4 GHz).  

Note that these are organized on a per channel basis, but sidebands are not (usually) considered separate channels, unless you're trying to get tricky with sideband averaging activities, which might be useful for concept instruments or other sensitivity studies. 

Here's an example what your \verb|SRFDATA/| directory could look like:
\begin{verbatim}
tropics_sv2_ch1.txt
tropics_sv2_ch2.txt
tropics_sv2_ch3.txt
tropics_sv2_ch4.txt
tropics_sv2_ch5.txt
tropics_sv2_ch6.txt
tropics_sv2_ch7.txt
tropics_sv2_ch8.txt
tropics_sv2_ch9.txt
tropics_sv2_ch10.txt
tropics_sv2_ch11.txt
tropics_sv2_ch12.txt
\end{verbatim}
filename structure is: \verb|sensor_id_ch##.txt|.  Not sure if ``Ch01'' works vs. ``Ch1''.  

\subsection{Creating a Case Directory}
The coefficient generation should be carried out in the \verb|workflows| folder.
In the case of a MW sensor, the \verb|workflows/MW| directory is the appropriate location.

In order to create a new case directory, simple type:
\begin{verbatim}
 ./Create_CaseDirectory -i <instrument_name>
\end{verbatim}
where \verb|instrument_name| is the CRTM instrument identifier. This will create the case directory and automatically link in all necessary executables.\\
For the concrete example of TROPICS this would be:
\begin{verbatim}
 ./Create_CaseDirectory -i tropics_sv1_srf_v2
\end{verbatim}
For help on the input options type:
\begin{verbatim}
 ./Create_CaseDirectory -h
\end{verbatim}
Once the case folder is created, you can enter it to proceed.
\begin{verbatim}
 cd tropics_sv1_srf_v2/
 ln -s <path_to_build>/bin/* .  
\end{verbatim}
This softlinks the previously compiled executables to the local working directory, you can also play around with your \verb|PATH| variable to point to the \verb|bin| directory. 

Note: When running these later, if you get an error about not being able to find a particular \verb|.so| library, you can add the paths to the required libraries using the \verb|LD_LIBRARY_PATH| variable, or be lazy (like me) and soft link the appropriate libraries to the local directory.   Please contact the JCSDA CRTM team for recommendations on better ways to handle this.   Building static isn't always possible, given that we depend on shared objects built on HPC stacks.

\subsection{Creating the oSRF file}
The CRTM oSRF file is an intermediate representation of the SRF as a netCDF file. The converter application is \verb|oSRF/oSRF_Create_from_ASCII|.
In order to obtain an oSRF file for TROPICS with 1 band, 12 channels and an interpolated SRF resolution of 1 MHz, follow the prompt:
\begin{verbatim}
> ./oSRF_Create_from_ASCII
 Enter the Sensor ID of the current instrument: 
tropics_sv1_srf_v2
 Enter the number of bands for the instrument: 
1
 Enter the number of instrument channels: 
12
 Is SRF interpolation requested (yes/no)?
yes 
 Enter the spectral resolution in GHz: 
0.001
\end{verbatim}
This will create the file \verb|tropics_sv1_srf_v2.osrf.nc|.  You can try it without interpolation too (recommended, in fact).  

\subsection{Creating the SpcCoeff file}
Creating the SpcCoeff file requires the \verb|polangle.txt| and \verb|polarization.txt| files to be linked into the case directory from the \verb|staging/fix/config/| directory.
They specify the polarization offset angle in degrees per channel and the CRTM polarization scheme per channel as an integer.
In order to create the SpcCoeff spectral coefficient file, run the \verb|Create_SpcCoeff| application in the case directory:

\verb|polarization.txt| structure, one entry per channel:
\begin{verbatim}
9
9
10
10
10
10
10
10
\end{verbatim}
Where these numbers refer to the values defined in:
\verb|./InstrumentInfo/SensorInfo/SensorInfo_Parameters.f90|
and repeated here for your convenience, but double check:
\begin{verbatim}
  INTEGER, PARAMETER :: N_POLARIZATION_TYPES    = 14
  INTEGER, PARAMETER :: INVALID_POLARIZATION    = 0
  INTEGER, PARAMETER :: UNPOLARIZED             = 1
  INTEGER, PARAMETER :: INTENSITY               = UNPOLARIZED
  INTEGER, PARAMETER :: FIRST_STOKES_COMPONENT  = UNPOLARIZED
  INTEGER, PARAMETER :: SECOND_STOKES_COMPONENT = 2
  INTEGER, PARAMETER :: THIRD_STOKES_COMPONENT  = 3
  INTEGER, PARAMETER :: FOURTH_STOKES_COMPONENT = 4
  INTEGER, PARAMETER :: VL_POLARIZATION         = 5
  INTEGER, PARAMETER :: HL_POLARIZATION         = 6
  INTEGER, PARAMETER :: plus45L_POLARIZATION    = 7
  INTEGER, PARAMETER :: minus45L_POLARIZATION   = 8
  INTEGER, PARAMETER :: VL_MIXED_POLARIZATION   = 9
  INTEGER, PARAMETER :: HL_MIXED_POLARIZATION   = 10
  INTEGER, PARAMETER :: RC_POLARIZATION         = 11
  INTEGER, PARAMETER :: LC_POLARIZATION         = 12
  INTEGER, PARAMETER :: CONST_MIXED_POLARIZATION= 13
  INTEGER, PARAMETER :: PRA_POLARIZATION        = 14
\end{verbatim}
Most passive MW sounders are 9/10.   Most passive MW imagers are 5/6.   If you're not sure, or if it's something weird, then discuss with the CRTM team(s).   

\verb|polangle.txt| structure, which defines the rotation of the polarization angle with respect to the along-flight direction, 0 degrees being directly along the direction of flight, assuming a standard sensor orientation (lots of caveats apply).  Usually these are given in degrees of offset clockwise from zero.   One entry per channel. 
\begin{verbatim}
0.00
0.00
<...> 
0.00
0.00
\end{verbatim}
Note: If you find you've got some biases WRT obs in your cross-track MW instrument over ocean, particularly at the edges of the scan, you might try some different polarization options. 

After you've created these files in your staging directory, copy the following to your working directory (e.g., \verb|workflows/MW/sensor|)
\begin{verbatim}
SRFDATA/
polangle.txt
polarization.txt
<sensor_id>_srf_v?.osrf.nc
\end{verbatim}

is what you should have so far in that directory (and the links to the binary files).

Yay.  Now we can finally create the Spectral Coefficient (aka \verb|SpcCoeff|) file.   
\begin{verbatim}
> ./Create_SpcCoeff 

     **********************************************************
                           Create_SpcCoeff

      Program to create the SpcCoeff files from the sensor 
      oSRF data files.

      $Revision: 46293 $
     **********************************************************


     Enter the oSRF netCDF filename: tropics_sv1_srf_v2.osrf.nc

     Default SpcCoeff version is: 1. Enter value: 2

\end{verbatim}
This will create the file \verb|tropics_sv1_srf_v2.SpcCoeff.nc|.

\subsection{Converting the SpcCoeff file from netCDF to Binary format}
Converting the SpcCoeff file from netCDF to binary format (if you want to) works like this:
\begin{verbatim}
 > ./SpcCoeff_NC2BIN 

     **********************************************************
                           SpcCoeff_NC2BIN

      Program to convert a CRTM SpcCoeff data file from netCDF 
      to Binary format.

      $Revision$
     **********************************************************


     Enter the INPUT netCDF SpcCoeff filename : tropics_sv1_srf_v2.SpcCoeff.nc

     Enter the OUTPUT Binary SpcCoeff filename: tropics_sv1_srf_v2.SpcCoeff.bin

     Increment the OUTPUT version number? [y/n]: n
\end{verbatim}
This will create the file \verb|tropics_sv1_srf_v2.SpcCoeff.bin| (either little-endian or big-endian, depending on your system).
If you need both endian formats of binary files, you'll have to recompile/rerun the \verb|SpcCoeff_NC2BIN| in the other endian.  e.g., \verb|-convert big_endian| flag.   We leave this as an exercise for the reader.

\subsection{Rosenkranz MW Transmittance Calculations}
The simplest approach to compute the line-by-line MW layer-to-space transmittances for the ODPS regression is to use the Rosenkranz model implemented in the CRTM.
First, link in the \verb|ECMWF83.AtmProfile.nc| predictor dataset from the \verb|data/reference/TauProd/| folder.
Second, create a \verb|SensorInfo_<sensor_id>| file by modifying an existing template or create a script to write the SensorInfo file for more complex instruments.
A generic SensorInfo generator will be provided in the future. 

\begin{verbatim}
 ./MW_TauProfile 

     **********************************************************
                            MW_TauProfile

      Program to compute transmittance profiles for 
      user-defined microwave instruments.

      $Revision: 45710 $
     **********************************************************


     Select atmospheric path
          1) upwelling  
          2) downwelling
     Enter choice: 1

     Enter the SensorInfo filename: SensorInfo_tropics
       Number of sensor entries read: 1

     Enter the netCDF AtmProfile filename: ECMWF83.AtmProfile.nc
\end{verbatim}

Depending on the spectral resolution of the oSRF file, these calculations may take several minutes.  

At this point your workflows directory should look something like:
\begin{verbatim}
 SRFDATA/
 polangle.txt
 polarization.txt
 ECMWF83.AtmProfile.nc
 <sensor_id>_srf_v?.osrf.nc
 <sensor_id>.SpcCoeff.nc
 <sensor_id>.SpcCoeff.bin
 upwelling.<sensor_id>.TauProfile.nc
\end{verbatim}

\subsection{ODPS Regression}
The ODPS regression is the last step in the CRTM MW transmittance coefficient generation and produces the TauCoeff file.
This step is performed in a separate folder and an example case directory for the regression is provided in \verb|workflows/MW/GenCoeff/|.
Before starting the regression calculations, the instrument ID needs to be added to a \verb|sensor_list| file that goes right in your workflows directory like everything else. 
\begin{verbatim}
 ## This file contains a list of sensors and is one of the input file of
 ## run_tau_coeff.sh
 ## Only those between BEGIN and END are processed.

 BEGIN_LIST
 tropics_sv1_srf_v2
 END_LIST
\end{verbatim}

Afterwards, the \verb|tau_coeff.parameters| input file needs to be modified for your local system.  I recommend copying from a previous example and modifying to suit your requirements.
We fully realize this isn't fun or very generic.  Future releases will include utilities that automate this process.  Here's an example from the \verb|cosmir_air| example I did on S4:
\begin{verbatim}
:MAX_CPUs:256
:CH_INT:1
:EXE_FILE:/data/users/bjohnson/CRTM/tmp3/CRTM_coef/build/bin/ODPS_Regress
:WORK_DIR:/data/users/bjohnson/CRTM/tmp3/CRTM_coef/workflows/MW/cosmir_air
:PROF_SET:ECMWF83
:SPC_COEFF_DIR:/data/users/bjohnson/CRTM/tmp3/CRTM_coef/workflows/MW/cosmir_air
:TAU_PROFILE_DIR:/data/users/bjohnson/CRTM/tmp3/CRTM_coef/workflows/MW/cosmir_air
:ATM_PROFLE_FILE:/data/users/bjohnson/CRTM/tmp3/CRTM_coef/workflows/MW/cosmir_air/ECMWF83.AtmProfile.nc
:GET_SEN_INFO:/data/users/bjohnson/CRTM/tmp3/CRTM_coef/build/bin/GetSenInfo
:IR_TYPE:3
:COMPONENTS:DEFAULT
#COMPONENTS:ozo
# Parameters below should not be modified in normal operations
:COMPONENT_GROUP1:dry,wlo,wco,ozo,co2,n2o,co,ch4
:COMPONENT_GROUP2:dry,wlo,wco,ozo,co2
:COMPONENT_GROUP3:dry,wet
\end{verbatim}
Make sure that all variables in \verb|tau_coeff.parameters| point to the right directories.
If the \verb|control/| doesn't exist, create it now:
\begin{verbatim}
 mkdir control
\end{verbatim}
Lastly, open the \verb|gen_tau_coeff.sh| script, scroll down, and check that the log output of the slurm script template points to the \verb|control/| directory:
\begin{verbatim}
echo "#!/bin/bash" >> ${jobScript}
echo "#SBATCH --job-name=ODPS_Run.%j.out" >> ${jobScript}
echo "#SBATCH --partition=serial" >> ${jobScript}
echo "#SBATCH --export=ALL" >> ${jobScript}
echo "#SBATCH --time=1:00:00" >> ${jobScript}
echo "#SBATCH --ntasks=1" >> ${jobScript}
echo "#SBATCH --cpus-per-task=1" >> ${jobScript}
echo "#SBATCH --mem-per-cpu=4096" >> ${jobScript}
echo "#SBATCH --output=<path>/CRTM_coef/workflows/MW/cosmir_air/control/output.%j.txt" >> ${jobScript}
echo "#SBATCH --error=<path>/CRTM_coef/workflows/MW/cosmir_air/control/output.%j.err" >> ${jobScript}
#echo "source /etc/bashrc" >> ${jobScript}
# Load the required modules (you might have to mess around with these)
echo "module purge" >> ${jobScript}
echo "module load license_intel/S4" >> ${jobScript}
echo "module load stack-intel/2021.5.0" >> ${jobScript}
echo "module load openblas" >> ${jobScript}
# Go to the execution directory
echo "cd ${RunDir}" >> ${jobScript}
#echo "I_MPI_JOB_STARTUP_TIMEOUT=10000" >> ${jobScript}
echo "${EXE_FILE} ${GROUP_ID}" >> ${jobScript}

# -- execute the job script
chmod 700 $jobScript
sbatch ./$jobScript
\end{verbatim}
You may wish to make substantial modifications to the above depening on your HPC environment.  This is also where you would do some modifications to run locally instead, not officially supported yet.

Finally, to perform the regression:


\begin{verbatim}
./gen_tau_coeff.sh tau_coeff.parameters
./cat_taucoef.sh tau_coeff.parameters ./Assemble_ODPS 
\end{verbatim}

If everything worked, the resulting TauCoeff file can be found in the \verb|results/| folder.  If it seems to hang, check your \verb|control/| directory for the output / errors.  Usually the issue is some library missing.   Feel free to ask for help.

You'll want to rename the resultant \verb|results/taucoef.<sensor_id>.ECMWF83.nc| file to \verb|<sensor_id>.TauCoeff.nc| to be consistent with typical CRTM naming convention.  
To convert to binary format, use \verb|ODPS_NC2BIN| and follow the prompts, with the same endian caveats as previously noted. 

\end{document}
